\section{Thách thức \& Giải pháp}

UDP không đảm bảo thứ tự, độ tin cậy hay kiểm soát tắc nghẽn, nên mọi vấn đề này phải được giải quyết ở tầng ứng dụng. Các thách thức chính và cách xử lý:

\begin{enumerate}[leftmargin=*]
    \item \textbf{Phân mảnh gói tin.} Datagram lớn hơn MTU bị phân mảnh ở tầng IP; mất một mảnh là mất cả datagram. \emph{Giải pháp}: giới hạn UDP payload 1400 byte, trừ header 22 byte và thẻ GCM 16 byte nên mỗi gói mang tối đa \textbf{1362 byte} dữ liệu. \texttt{Stream} tự cắt dữ liệu lớn thành chunk, gắn \texttt{offset}, chunk cuối mang \texttt{FLAG\_EOM}; bên nhận ghép lại theo \texttt{offset} (chịu được đảo thứ tự) với giới hạn bộ đệm 16~MB để chống cạn bộ nhớ.

    \item \textbf{Mất gói và đảo thứ tự.} \emph{Giải pháp}: ACK chọn lọc (SACK) cho mỗi gói \texttt{DATA}; RTO tính theo RTT (RFC~6298, kẹp 0,2--10~s); fast retransmit khi có $\ge 3$ gói số cao hơn đã được ACK; tối đa 8 lần truyền lại. Chỉ gửi ACK \emph{sau khi} giải mã thành công để gói giả mạo không làm sai trạng thái.

    \item \textbf{Tắc nghẽn và tràn bộ đệm.} Gửi file bằng tốc độ tối đa làm mất gói hàng loạt. \emph{Giải pháp}: cửa sổ tắc nghẽn CUBIC ($C=0{,}4$, $\beta=0{,}7$, cwnd khởi đầu 10 gói); hàm \texttt{can\_send()} điều tiết hàng đợi gửi.

    \item \textbf{Deadlock giữa luồng nhận và luồng gửi.} Server dùng một luồng cho cả nhận và gửi; nếu hàm gửi \emph{chặn chờ} cửa sổ tắc nghẽn mở, thì luồng đó không thể xử lý ACK để mở cửa sổ $\Rightarrow$ treo vĩnh viễn. \emph{Giải pháp}: \texttt{\_flush\_send\_queue()} \textbf{không bao giờ chặn}: dữ liệu vào hàng đợi \texttt{\_tx\_queue}, được xả khi (i) vừa enqueue, (ii) vừa xử lý ACK, (iii) mỗi \texttt{tick()}. Ngoài ra dùng \texttt{RLock} cho khóa gửi (cho phép gọi lồng nhau) và tách khóa riêng cho số thứ tự và bảng stream, không bao giờ giữ hai khóa theo thứ tự đảo.

    \item \textbf{Mất gói bắt tay.} \texttt{HELLO}/\texttt{HELLO\_ACK} không có ACK riêng. \emph{Giải pháp}: client gửi lại \texttt{HELLO} mỗi 1~s (tối đa 5 lần), server lưu \texttt{HELLO\_ACK} và phát lại khi nhận \texttt{HELLO} trùng.

    \item \textbf{Đổi địa chỉ/cổng (NAT rebinding) và giả mạo.} Định danh theo \texttt{conn\_id} cho phép đổi mạng, nhưng kẻ tấn công có thể giả IP. \emph{Giải pháp}: chỉ cập nhật địa chỉ khi datagram từ địa chỉ mới là \texttt{DATA} mã hóa giải mã được bằng khóa phiên cũ (\texttt{verify\_migration\_datagram}).

    \item \textbf{Tái sử dụng nonce AES-GCM.} Dùng lại nonce phá vỡ an toàn của GCM. \emph{Giải pháp}: khóa riêng cho từng chiều, nonce dựng từ \texttt{seq\_num}, kèm bộ kiểm tra nonce đã dùng; header được ràng buộc làm AAD nên sửa header sẽ làm giải mã thất bại.

    \item \textbf{Kết nối ``chết'' sau NAT/firewall.} \emph{Giải pháp}: heartbeat 5~s; sau 30~s không nhận gì thì đóng phiên và giải phóng tài nguyên, thông báo lại danh sách người dùng.
\end{enumerate}

\textbf{Bài học.} Việc tách rõ các tầng (gói tin $\rightarrow$ độ tin cậy $\rightarrow$ stream $\rightarrow$ ứng dụng) giúp kiểm thử từng phần bằng các bộ test riêng (\texttt{tests/}) và tránh các lỗi đồng bộ khó phát hiện.
